\section{Calibrating the physical rotation action}
\label{sec:physical-angular-actions}

We construct the rotation action on the physical field spaces from the
actual angular Ward maps. Equal dimensions of the physical and reference
spaces would not supply this action. The essential step is to correct an
approximate intertwiner while preserving the real structure and signed
coordinate symmetries.

Write $E_d$ and $E_d^0$ for the finite-dimensional complex physical and
reference field spaces at grade cap $d$. Retain the actual derivative maps $\partial$, the conserved
stress $J$, its translation charges $p$, and its angular maps $Y$.
The physical Ward propagation argument supplies one real $c\ne0$ such
that $p=c\partial$ at every cap. Put $D=c^{-1}Y$. Its output also
supplies the global angular Ward identity for $D$, the Lie-bracket
identity
\[
                 [D_d(A),D_d(B)]=D_d([A,B]),
\]
and compatibility with conjugation, signed motions and inclusions.
These are the inputs from the physical shell argument, distinct from the
reference Maxwell Ward calculation.

On the reference spaces keep the fixed continuous rotation representation
$R_d^0$ and its infinitesimal maps $C_d$. Its reference stress has Ward
coefficient $\kappa\ne0$, with the fixed value $-1$ in
Section~\ref{sec:free-maxwell-ward-flux}. The reference action and its
Lorentz companion are the ones attached to that same free hierarchy.
We do not replace either representation when calibrating the physical one.

\subsection{The scalar normalization on a common comparison sequence}

Fix $d$. The finite-request comparison gives, on one sequence, complex-linear isomorphisms
$U_{n,e}:E_e^0\to E_e$, radii $r_n>0$ and stress normalizers $\xi_n>0$.
Choose it sufficiently rich to include the derivative test at cap $4$,
the translation-charge test there, the angular tests at cap $d$, and
the real/signed compatibility at $d$. One can do this by taking the
maximum of the finitely many requested caps, arities, derivative orders
and compact radii before invoking the comparison. This is the common
finite-request construction of
Section~\ref{sec:actual-stress-normalization}; the shell-charge limits
are supplied by the physical Ward comparison.

Transport both source and target of the derivative maps through these
same $U_{n,e}$. Denote the transported cap-$4$ derivative array,
multiplied by $r_n$, by $A_n$. The two limits just specified say
\[
        A_n\longrightarrow A^0\ne0,\qquad
        (c/\xi_n)A_n\longrightarrow\kappa A^0.
\]
The nonzero derivative array is the fixed reference test. Choose a real
linear functional $\lambda$ on the finite-dimensional array space with
$\lambda(A^0)\ne0$. Then $\lambda(A_n)\ne0$ eventually, and division
gives
\begin{equation}\label{eq:angular-scalar-ratio}
 \frac c{\xi_n}
  =\frac{\lambda((c/\xi_n)A_n)}{\lambda(A_n)}
                         \longrightarrow\kappa.
\end{equation}
In particular no unproved lower bound on $\xi_n$ is required.
The angular comparison is
\[
 \xi_n^{-1}U_{n,d}^{-1}Y_d(A)U_{n,d}
                             \longrightarrow\kappa C_d(A).
\]
Multiplying by $(c/\xi_n)^{-1}$ and using
\eqref{eq:angular-scalar-ratio} proves
\begin{equation}\label{eq:angular-transport-limit}
 B_n(A):=U_{n,d}^{-1}D_d(A)U_{n,d}\longrightarrow C_d(A).
\end{equation}
Only finitely many generating planes are needed at this stage. The
maps $U_{n,d}$ commute with conjugation and the signed actions.
Their convergence is neither required nor asserted.

\subsection{An exact intertwiner from nearby infinitesimal matrices}

We give the finite-dimensional argument used with
\eqref{eq:angular-transport-limit}. Let $E=E_d^0$ and write
$K_{ab}=e_ae_b^{\mathsf T}-e_be_a^{\mathsf T}$ for the ordered plane
generators, including $K_{aa}=0$. Index them by $i=(a,b)$. For every
pair choose an axis $h$ distinct from $a,b$. Direct multiplication gives
\[
 [K_{ah},K_{hb}]=K_{ab}.
\]
Thus the corresponding matrices $B_{n,i}$ and $C_i$ obey a fixed family
of perfect-commutator identities $B_{n,i}=[B_{n,l(i)},B_{n,r(i)}]$,
and similarly for $C_i$. On $\operatorname{End}_{\mathbb C}(E)$ put
\[
 T_{n,i}(F)=B_{n,i}F-FC_i,\qquad T_i(F)=C_iF-FC_i,
 \qquad L_n=\sum_iT_{n,i}^2,\quad L=\sum_iT_i^2.
\]
Left and right multiplication commute, so the $T_{n,i}$ satisfy the
same perfect identities. In fact the full Lie-bracket relations give
\begin{equation}\label{eq:angular-casimir-commutes}
                        [L_n,T_{n,i}]=0.
\end{equation}
For completeness, the plane bracket formula is
\[
 [K_{ab},K_{ij}]=\delta_{bi}K_{aj}-\delta_{bj}K_{ai}
                      +\delta_{ai}K_{jb}-\delta_{aj}K_{ib}.
\]
Insert this in
$[T_{ab}^2,T_{ij}]=T_{ab}[T_{ab},T_{ij}]+[T_{ab},T_{ij}]T_{ab}$
and sum over the ordered pairs $(a,b)$. The four resulting sums cancel
in pairs after exchanging their repeated indices and using
$T_{ba}=-T_{ab}$. This proves \eqref{eq:angular-casimir-commutes}
for the relative representation too.

First construct a real polynomial $q$ such that $P=q(L)$ is a projection,
$P(I)=I$, and $T_iP=0$ for every $i$. Average any positive-definite real
inner product on $\operatorname{End}_{\mathbb C}(E)$ under
$F\mapsto R_d^0(K)F R_d^0(K)^{-1}$, using probability Haar measure on
the compact group $O(4)$. The average is positive definite: for a
nonzero vector the integrand is positive at the identity and, by
continuity, on a nonempty open set of positive Haar measure. It is
invariant by Haar translation. Differentiating that invariance makes
each $T_i$ skew for the averaged form $\langle\cdot,\cdot\rangle_0$.
Consequently
\[
       \langle LF,F\rangle_0=-\sum_i\langle T_iF,T_iF\rangle_0,
       \qquad \ker L=\bigcap_i\ker T_i.
\]
Moreover $L^2F=0$ implies $LF=0$: the displayed kernel identity gives
$T_iLF=0$, and skewness then gives
$\langle LF,LF\rangle_0=\sum_i\langle T_i^2F,LF\rangle_0=0$.

Cayley--Hamilton provides a nonzero polynomial annihilating $L$.
Starting with $Lp(L)=0$, divide out factors of the variable as long as
$p(0)=0$. If $p(X)=Xh(X)$, then $L^2h(L)=0$ implies $Lh(L)=0$ by the
preceding implication, applied to every vector. Eventually a polynomial
with nonzero constant term remains. Normalize it to obtain
$q(0)=1$ and $Lq(L)=0$. Since $q(L)$ is the identity on $\ker L$,
it is a projection onto that kernel. The identity matrix belongs to
this kernel, so $P(I)=I$. This proof also covers $E=\{0\}$.

Let $\mathcal F$ be the closed real vector space of endomorphisms of $E$
commuting with the reference signed actions and conjugation. The identity
belongs to $\mathcal F$. Every $L_n$ preserves $\mathcal F$.
Indeed, a signed coordinate transformation permutes the ordered planes
and multiplies each by a sign. Both $B_{n,i}$ and $C_i$ acquire the same
sign, so it cancels in $T_{n,i}^2$ and the sum is unchanged. Conjugation
commutes with both matrix families; since the polynomial coefficients
are real, it commutes with all subsequent polynomial operations too.
It follows that $P_n=q(L_n)$ preserves $\mathcal F$, commutes with
every $T_{n,i}$, and converges to $P$.

The $P_n$ need not be projections. Correct them by the polynomial
iteration
\[
       Z_{j+1}=3Z_j^2-2Z_j^3,\qquad Z_0=P_n.
\]
For $e(Z)=Z^2-Z$ the exact identities are
\[
 Z_{j+1}-Z_j=(I-2Z_j)e(Z_j),\qquad
 e(Z_{j+1})=(4e(Z_j)-3I)e(Z_j)^2.
\]
Here is a uniform convergence estimate. Choose $A>\|P\|$ and
$C>\|I\|+2A$. For sufficiently large $n$, put
$\varepsilon_n=\|e(P_n)\|$; then
\[
 \|P_n\|+2C\varepsilon_n\leq A,\qquad
 (4\varepsilon_n+3\|I\|)\varepsilon_n\leq\tfrac12.
\]
Induction using the two identities gives
$\|e(Z_j)\|\leq2^{-j}\varepsilon_n$ and
$\|Z_{j+1}-Z_j\|\leq C2^{-j}\varepsilon_n$.
Their sum keeps $\|Z_j\|\leq A$, closing the induction. Thus
$Z_j\to Q_n$, with $Q_n^2=Q_n$ and
$\|Q_n-P_n\|\leq2C\varepsilon_n$. In particular $Q_n\to P$.
Closedness preserves the constraints: $Q_n$ preserves $\mathcal F$
and commutes with all $T_{n,i}$.

The restricted operators $T_{n,i}Q_n$ satisfy the same perfect identities,
since $Q_n$ is idempotent and commutes with them. They converge to
$T_iP=0$. Let $m_n=\max_i\|T_{n,i}Q_n\|$. The perfect identities give
$m_n\leq2m_n^2$. Eventually $m_n<1/2$, hence $m_n=0$. Therefore
\[
                        T_{n,i}Q_n=0\quad\text{for every }i.
\]
Apply $Q_n$ to the identity matrix and set $F_n=Q_n(I)$. Then
$F_n\in\mathcal F$, $F_n\to I$, and
$B_{n,i}F_n=F_nC_i$. For one sufficiently late $n$, $F_n$ is invertible
by the Neumann series. Thus
\begin{equation}\label{eq:angular-exact-intertwiner}
              Q_d=U_{n,d}F_n:E_d^0\longrightarrow E_d
\end{equation}
is an isomorphism preserving reality and signed motions, and
$D_d(A)Q_d=Q_dC_d(A)$ for every skew matrix $A$.
To pass from planes to all $A$, use the exact spanning identity
$A=\frac12\sum_{a,b}A_{ab}K_{ab}$. No complete reducibility or
unitarity of the deformed matrices was assumed.

\subsection{The physical action and its compatibility}

Choose one map \eqref{eq:angular-exact-intertwiner} in each cap and define
\[
                         R_d(K)=Q_dR_d^0(K)Q_d^{-1}.
\]
It is a continuous representation of all of $O(4)$, preserves the real
structure, agrees with the prescribed signed action, and has derivative
$D_d(A)$ along $\exp(tA)$ at zero. The group law moves this derivative
to every $t$. The choices of $Q_d$ need not commute with inclusions.

For arbitrary distinct points and field labels, differentiate
\[
 t\longmapsto S(R_{d_1}(e^{tA})v_1,\ldots,
                     R_{d_n}(e^{tA})v_n;e^{tA}x_1,\ldots,e^{tA}x_n).
\]
Joint differentiability follows by expanding the labels in finite bases:
the expression is a finite sum of smooth position coefficients times
multilinear coordinate products. The chain rule gives the global angular
Ward expression for $D$, which is zero. Thus the function is constant.
Simultaneous covariance holds for every one-parameter rotation and also
for signed transformations. The set of transformations for which it
holds is a subgroup, by composition and inversion. Plane rotations
generate $SO(4)$, and a signed reflection supplies its other component,
as in the group-generation argument of
Section~\ref{sec:lorentz-covariance}. Hence covariance holds for all $O(4)$.

The unit, inclusions and derivatives are now compatible with this actual
action. For the unit, test $R_d(K)1$ against any separated spectator word,
move all labels and positions back by $K^{-1}$, and delete the unit using
the original unit law. Move the remaining word forward again. The result
is the test of $1$, so faithfulness of the physical quotient gives
$R_d(K)1=1$. For an inclusion $\iota_{de}$, the identical argument uses
the original inclusion identity at the moved configuration and yields
\[
                     R_e(K)\iota_{de}=\iota_{de}R_d(K).
\]
Finally differentiate the covariance identity in the coordinate of the
head insertion, keeping arbitrary spectators. The ordinary linear chain
rule and the original derivative identity give
\begin{equation}\label{eq:angular-derivative-vector}
 R_{d+1}(K)\partial_a v
                  =\sum_i K_{ia}\partial_i(R_d(K)v).
\end{equation}
Again all-spectator faithfulness converts equality of correlations into
an identity of field classes. This proves the full physical Euclidean
action laws rather than merely laws at each isolated cap.

\subsection{Normalizing and averaging the stress}

First replace $J$ by $J_0=c^{-1}J$. Linearity of the collar flux makes
its translation charge exactly $\partial$ and its angular charge exactly
$D$. The tensor remains real, symmetric and divergence-free. To make it
transform tensorially under the action just constructed, define
\[
 (\mathcal R_KJ_0)_{ab}
               =\sum_{i,j}K_{ai}K_{bj}R_4(K)(J_{0,ij}),
 \qquad \widehat J=\int_{O(4)}\mathcal R_KJ_0\,dK.
\]
This integral is in the finite-dimensional space $E_4^{16}$, with
probability Haar measure. Its integrand is continuous. The group law
gives $\mathcal R_K\mathcal R_L=\mathcal R_{KL}$, so Haar invariance
implies $\mathcal R_K\widehat J=\widehat J$ and hence
\[
        R_4(K)\widehat J_{ij}
                    =\sum_{a,b}K_{ai}K_{bj}\widehat J_{ab}.
\]
The average is real and symmetric. Equation~\eqref{eq:angular-derivative-vector}
and orthogonality show
\[
 \sum_a\partial_a(\mathcal R_KJ_0)_{ab}
       =\sum_jK_{bj}R_5(K)\left(\sum_i\partial_iJ_{0,ij}\right)=0.
\]
The derivative maps are continuous linear maps of finite-dimensional
spaces and therefore commute with the integral. Thus $\widehat J$ is
conserved.

We verify that averaging retains both exact Ward normalizations.
For vectors $u,v$ write $J[u,v]=\sum_{a,b}u_av_bJ_{ab}$. Then
\[
 (\mathcal R_KJ)[u,v]=R_4(K)J[K^{-1}u,K^{-1}v].
\]
In a collar flux with cutoff $\chi$ and vector field $W$, change variables
$y=Kz$, moving the external points and labels by $K^{-1}$. Orthogonal
maps preserve volume, and
$\nabla(\chi\circ K)=K^{-1}(\nabla\chi)\circ K$.
The transformed vector field is $K^{-1}W(Kz)$. Thus the flux of
$\mathcal R_KJ_0$ equals the flux of $J_0$ in this transformed collar.
The compact support of $\nabla\chi$ avoids all insertion points, so the
integrands are smooth there and the change of variables is legitimate.

For a translation $W=e_a$, the transformed charge is the linear
combination with coefficients $(K^{-1}e_a)_i$. Its Ward identity,
\eqref{eq:angular-derivative-vector} and covariance of the external
correlation return exactly the original derivative charge $\partial_a$.
For $W(y)=A(y-x_0)$ the transformed field is
$(K^{-1}AK)(z-K^{-1}x_0)$. Differentiating the representation identity
$R_d(K)R_d(e^{tA})R_d(K)^{-1}=R_d(e^{tKAK^{-1}})$ gives
\[
               R_d(K)D_d(A)=D_d(KAK^{-1})R_d(K).
\]
Its use in the transformed Ward identity returns exactly $D_d(A)$.
This also holds for orientation-reversing $K$; volume measure contributes
$|\det K|=1$, not a determinant sign. Each fixed collar flux is a
continuous linear functional of the stress labels. It therefore passes
through the finite-dimensional Haar integral, and the constant target
Ward expression integrates to itself. Both charges of $\widehat J$
are consequently normalized as required.

\subsection{Keeping the same maps at the Lorentz step}

The maps $Q_d$ in \eqref{eq:angular-exact-intertwiner}, with the reference
Euclidean and Lorentz actions already fixed, give precisely the calibrated
actions of Input~\ref{input:calibrated-actions}. Their reality and signed
compatibility were preserved by the polynomial correction; their
Euclidean actions now commute with inclusions. Define the physical
Lorentz action at each cap by $Q_dT_d^0Q_d^{-1}$ with these very maps.
Proposition~\ref{prop:lorentz-label-descent} proves its cross-cap
compatibility from the boost differential equation, and
Section~\ref{sec:lorentz-covariance} supplies the boundary covariance
argument. There is no new choice of a Lorentz intertwiner. The hierarchy,
its derivative maps and its field quotients have remained the same
through this construction.
